\chapter{Пространственное связывание составной хиральной монеты}
\label{ch:v6-spectral-transition-discrete-composite-higgs-spatial-binding-gate}

\section{Минимальный пространственный закон}

Предыдущий гейт построил эндогенный дублет
$H_{\rm eff}=\ell\overline e$, но не проверял локализацию. Простая монета,
действующая только на $L_L\oplus e_R$ и не смешивающая направления,
структурно не может связать пакет: нормы правого и левого потоков
сохраняются отдельно и каждый поток сдвигается на один узел.

Минимальный кандидат, способный развернуть поток без внешнего массового
угла, имеет генератор
\begin{equation}
 G_n(\Psi)=\sigma_y^{\rm dir}\otimes K(H_{{\rm eff},n}),
 \label{eq:v6-composite-spatial-generator}
\end{equation}
где
\begin{equation}
 H_{{\rm eff},n}
 =\sum_{d=\pm}\ell_{n,d}\overline{e_{n,d}}
\end{equation}
и $K(H)$ является самосопряжённым хиральным блоком предыдущего гейта.
Один шаг равен
\begin{equation}
 \Psi^{t+1}=S\exp[-i\kappa G(\Psi^t)]\Psi^t,
 \label{eq:v6-composite-spatial-update}
\end{equation}
где $S$ сдвигает два направления в противоположные стороны.

Правило ближайшее, точно унитарно на каждом узле и глобально калибровочно
ковариантно. В отличие от внешнего массового угла, весь внутренний
генератор исчезает вместе с состоянием.

\section{Линеаризация около вакуума}

Составной дублет квадратичен по амплитуде:
\begin{equation}
 H_{\rm eff}(\varepsilon\Psi)=O(\varepsilon^2).
\end{equation}
Поэтому действие монеты на состояние отличается от тождества только в
третьем порядке:
\begin{equation}
 \exp[-i\kappa G(\varepsilon\Psi)]\varepsilon\Psi
 =\varepsilon\Psi+O(\varepsilon^3).
 \label{eq:v6-composite-cubic-linearization}
\end{equation}
Машинный scaling-тест даёт показатель $2.99999\ldots$.

Следовательно, линейный оператор около пустого вакуума равен чистому
безмассовому сдвигу. У него нет спектральной щели, в которой могла бы
лежать малая экспоненциально локализованная мода. Это точный структурный
запрет на малую гапированную частицу, но не на конечноподдержанное
нелинейное решение конечной амплитуды.

\section{Проспективный численный тест}

Проверен один заранее фиксированный гауссов пакет ширины четыре на $256$
узлах. Для
\begin{equation}
 \kappa\in\{0,0.5,1,2,4,8,16,32\}
\end{equation}
выполнено $80$ шагов без массового угла, внешнего потенциала и жёсткости.
Во всех ветвях:
\begin{enumerate}
 \item второй момент вырос с $4$ более чем до $3000$;
 \item IPR стал меньше начального;
 \item вероятность в ядре $|x|\le8$ стала меньше $0.1$;
 \item норма сохранилась с ошибкой ниже $10^{-14}$.
\end{enumerate}
Большая $\kappa$ замедляет баллистическое расхождение второго момента, но
не создаёт связанного профиля. Одновременно плотность становится более
распределённой, а не концентрированной.

Этот скан является stop-сертификатом для объявленного канонического пакета,
но не доказательством отсутствия всех нелинейных решений. Из-за безмассовых
хвостов остаётся специальная возможность: профиль с точно конечной опорой,
на границе которой нелинейная монета полностью гасит исходящий поток.

\begin{theorem}[Гапированное связывание не возникает]
Минимальная составная пространственная монета
\eqref{eq:v6-composite-spatial-update} локальна, нормосохраняюща и
ковариантна. Однако её линеаризация около вакуума равна безмассовому
сдвигу и не имеет щели. Проспективный скан восьми значений $\kappa$ не
находит устойчивого связанного профиля. Поэтому малая экспоненциально
локализованная частица не выводится; открыты только конечноподдержанные
решения специальной амплитуды.
\end{theorem}

\begin{nogocriterion}
Нельзя интерпретировать уменьшение скорости разбегания при большой
$\kappa$ как массу или связывание. Для связанного endpoint нужны
ограниченный второй момент и ненулевая асимптотическая доля нормы в ядре.
Также нельзя исключать все солитоны одним временным сканом: отдельного
аналитического теста требует compact-support loophole.
\end{nogocriterion}

\section{Следующий гейт}

Следующий гейт
\path{version6_spectral_transition_discrete_compacton_existence_gate}
должен решить точные конечномерные уравнения для одно- и двухузлового
профиля, возвращающегося после шага с общей фазой и без утечки опоры.

Нулевая гипотеза: граничный сдвиг запрещает ненулевой компактон либо
оставляет только решения при настроенном произведении $\kappa\|\Psi\|^2$.
Стоп-критерий: если решение требует непрерывной настройки $\kappa$ или
нормы, оно не проходит R5 даже при точной локализации.

\begin{protocol}
\begin{enumerate}
 \item локальный составной пространственный закон: \textbf{построен};
 \item сохранение нормы и ковариантность: \textbf{точные};
 \item порядок нелинейной поправки около вакуума: \textbf{третий};
 \item линейная массовая щель: \textbf{отсутствует};
 \item малая гапированная связанная мода: \textbf{запрещена};
 \item робастная локализация в проспективном скане: \textbf{не найдена};
 \item конечноподдержанный compacton: \textbf{остаётся открыт};
 \item физический масштаб и $\kappa$: \textbf{не выведены};
 \item R4 и R5: \textbf{не закрыты}.
\end{enumerate}
\end{protocol}

Машинный сертификат сохранён в
\path{s2t_v6_spectral_transition_discrete_composite_higgs_spatial_binding_gate_results.json}.